\documentclass[a4paper,12pt]{article}
\usepackage{fontspec}
\usepackage[main=russian,english]{babel}
\setmainfont{DejaVu Serif}
\setsansfont{DejaVu Sans}
\setmonofont{DejaVu Sans Mono}
\usepackage[a4paper,margin=24mm]{geometry}
\usepackage{array,booktabs,multirow,siunitx,tabularx}
\usepackage{hyperref}
\sisetup{output-decimal-marker={,}}
\renewcommand{\arraystretch}{1.18}
\title{Лабораторная работа №5\\Таблицы в LaTeX}
\author{Тесленко Кирилл Дмитриевич}
\date{}
\begin{document}
\maketitle
\section{Простая таблица}
В окружении \texttt{tabular} спецификаторы \texttt{l}, \texttt{c} и \texttt{r}
задают выравнивание по левому краю, центру и правому краю.
\begin{table}[htbp]
\centering
\caption{Основные параметры демонстрационного набора данных}
\label{tab:parameters}
\begin{tabular}{@{}lcr@{}}
\toprule
Параметр & Обозначение & Значение\\
\midrule
Число записей & $N$ & 240\\
Число признаков & $p$ & 8\\
Число групп & $k$ & 3\\
\bottomrule
\end{tabular}
\end{table}
\section{Объединение столбцов}
Команда \verb|\multicolumn| позволяет объединить несколько ячеек одной строки.
\begin{table}[htbp]
\centering
\caption{Сводная таблица для трёх вариантов}
\label{tab:summary}
\begin{tabular}{@{}lcc@{}}
\toprule
\multicolumn{3}{c}{Учебные показатели (условный пример)}\\
\midrule
Вариант & Время, мс & Ошибки, шт.\\
\midrule
A & 12,50 & 2\\
B & 9,75 & 1\\
C & 14,20 & 3\\
\bottomrule
\end{tabular}
\end{table}
\section{Выравнивание числовых данных}
Пакет \texttt{siunitx} использует тип столбца \texttt{S}: числа
выравниваются по десятичному разделителю. В исходном коде
разделитель --- точка, в готовом документе --- запятая.
\begin{table}[htbp]
\centering
\caption{Числовые значения с выравниванием по запятой}
\label{tab:numbers}
\begin{tabular}{@{}l S[table-format=2.2]@{}}
\toprule
Серия & {Время, мс}\\
\midrule
A & 1.25\\
B & 12.50\\
C & 9.80\\
\bottomrule
\end{tabular}
\end{table}
\section{Объединение строк и перенос длинного текста}
Пакет \texttt{multirow} объединяет соседние строки, а \texttt{tabularx}
создаёт колонку \texttt{X}, занимающую оставшуюся ширину таблицы.
\begin{table}[htbp]
\centering
\caption{Сравнение способов оформления}
\label{tab:comparison}
\small
\begin{tabularx}{\textwidth}{@{}l l X@{}}
\toprule
Категория & Средство & Назначение\\
\midrule
\multirow{2}{*}{Структура} & \texttt{tabular} & Задание строк, столбцов и выравнивания обычного текста.\\
 & \texttt{multicolumn} & Объединение ячеек заголовка по горизонтали.\\
\midrule
\multirow{2}{*}{Оформление} & \texttt{booktabs} & Горизонтальные линии без тяжёлой сетки.\\
 & \texttt{siunitx} & Аккуратное выравнивание десятичных чисел.\\
\bottomrule
\end{tabularx}
\end{table}
\section{Перекрёстные ссылки}
Таблица~\ref{tab:parameters} показывает простое выравнивание столбцов,
а таблица~\ref{tab:numbers} --- выравнивание чисел по десятичному разделителю.
\section{Вывод}
Выполнены четыре примера. Проверены \texttt{tabular}, \texttt{booktabs},
\texttt{multicolumn}, \texttt{multirow}, \texttt{tabularx} и \texttt{siunitx}.
Все числовые значения приведены для демонстрации вёрстки,
а не как результаты измерений.
\end{document}
